%Chargement des paquets
\usepackage{amsmath}
\usepackage{amsfonts}
\usepackage{amssymb}
\usepackage{enumerate}
\usepackage{mathtools}
\usepackage{bbm}
\usepackage{xparse, etoolbox}
\usepackage{enumerate}
\usepackage{enumitem}
\usepackage{mathabx}
\usepackage{babel}[french]

%environment
\usepackage{amsthm}
\usepackage{keytheorems}
\usepackage{hyperref} 

%Interface théorème
\newkeytheorem{lem}[
	name = Lemme
]

\newkeytheorem{prop}[
	name = Proposition
]

\newkeytheorem{thm}[
	name = Théorème
]

\newkeytheorem{coro}[
	name = Corollaire
]

\renewcommand*{\proofname}{Démonstration}

\theoremstyle{definition}
\newtheorem{defn}{Définition}

\theoremstyle{definition}
\newtheorem*{nota}{Notation}

\theoremstyle{definition}
\newtheorem*{limi}{Liminaire}

\theoremstyle{remark}
\newtheorem{rmq}{Remarque}

\theoremstyle{plain}
\newtheorem*{fait}{Fait}


%Ensembles de nombres
\newcommand{\N} {\mathbb{N}}
\newcommand{\Ne}{\N^\ast}
\newcommand{\Z} {\mathbb{Z}}
\newcommand{\Q} {\mathbb{Q}}
\newcommand{\R} {\mathbb{R}}
\newcommand{\Rb}{\overline{\mathbb{R}}}
\newcommand{\Rp}{\R_+}
\newcommand{\Rm}{\R_-}
\newcommand{\K} {\mathbb{K}}
\newcommand{\Ket} {\mathbb{K^{\ast}}}
\newcommand{\Cx}{\mathbb{C}}

%Ensembles, tribus
\newcommand{\A}{\mathbb{A}}
\renewcommand{\P}{\mathcal{P}}
\newcommand{\T}{\mathcal{T}}
\newcommand{\F}{\mathcal{F}}
\newcommand{\C} {\mathcal{C}}
\newcommand{\ouv}{\mathcal{O}}
\newcommand{\B}{\mathcal{B}}
\newcommand{\M}{\mathcal{M}}
\newcommand{\etag}{\mathcal{E}}
\newcommand{\etagOp}{\etag^{+}(\Omega)}
\newcommand{\etagO}{\etag(\Omega)}
%%Lp avant quotientage
\newcommand{\Lic}{\mathcal{L}^1}
\newcommand{\Lpc}{\mathcal{L}^p}
\newcommand{\Linfc}{\mathcal{L}^{\infty}}
\newcommand{\Lifull}{\Lic(\Omega, \B, m)}
\newcommand{\Lpfull}{\Lpc(\Omega, \B, m)}
\newcommand{\Linffull}{\Linfc(\Omega, \B, m)}
%%Lp après quotientage
\newcommand{\Li}{L^1}
\newcommand{\Ld}{L^2}
\newcommand{\Lp}{L^p}
\newcommand{\Linf}{L^{\infty}}
\newcommand{\Listd}{\Li(\Omega, m)} 
\newcommand{\Ldstd}{\Ld(\Omega, m)} 
\newcommand{\Lpstd}{\Lp(\Omega, m)} 

%Quelques opérateurs
\DeclareMathOperator{\codim}{codim}
\DeclareMathOperator{\rg}{rg}
\DeclareMathOperator{\tr}{tr}
\DeclareMathOperator{\vect}{Vect}
\DeclareMathOperator{\ima}{Im}
\DeclareMathOperator{\id}{Id}
\DeclareMathOperator{\sign}{sign}
\DeclareMathOperator{\ext}{ext}
\newcommand{\ind}{\mathbbm{1}}
\newcommand*{\spr}[2]{\left(\,#1\,\middle|\,#2\,\right)}
\newcommand{\interior}[1]{%
  {\kern0pt#1}^{\mathrm{o}}%
}
\DeclareMathOperator{\card}{card}
\DeclareMathOperator{\ppcm}{ppcm}

%Commandes de pure flemme
\newcommand{\veps}{\varepsilon}
\newcommand{\intab}{\int_{a}^{b}}
\newcommand{\intR}{\int_{-\infty}^{\infty}}
\newcommand{\intinf}{\int_{- \infty}^{+ \infty}}
\newcommand{\equi}{\Leftrightarrow}
\newcommand{\Kev}{$\K$-espace vectoriel }
\newcommand{\Kevs}{$\K$-espaces vectoriels }
\newcommand{\Rev}{$\R$-espace vectoriel }
\newcommand{\Revs}{$\R$-espaces vectoriels }
\newcommand{\Cev}{$\Cx$-espace vectoriel }
\newcommand{\Cevs}{$\Cx$-espaces vectoriels }
\newcommand{\evn}{(E, \norm{.})}
\newcommand{\interf}[2]{[#1,#2]}
\newcommand{\limb}{\lim\limits_}
\newcommand{\lebn}{\lambda_n}
\newcommand{\pp}{\text{~presque partout}}
\newcommand{\derivp}[2]{\frac{\partial #1}{\partial #2}}
\newcommand{\famiu}{(u_i)_{i \in I}}
\newcommand{\sev}{\text{~s.e.v.}}
\newcommand{\Mnp}{M_{n,p}(\K)}
\newcommand{\Mn}{M_{n}(\K)}
\newcommand{\tq}{\text{~t.q.~}}
\newcommand{\mq}{~montrer~que~}
\newcommand{\et}{\quad \text{et} \quad}
\newcommand{\ou}{\text{~ou~}}
\newcommand{\Kx}{\K[X]}


%Commandes de gars chelou
\renewcommand{\subset}{\subseteq}

%Commande restriction, ty duckduckgo
\def\restr#1#2{\mathchoice
              {\setbox1\hbox{${\displaystyle #1}_{\scriptstyle #2}$}
              \restraux{#1}{#2}}
              {\setbox1\hbox{${\textstyle #1}_{\scriptstyle #2}$}
              \restraux{#1}{#2}}
              {\setbox1\hbox{${\scriptstyle #1}_{\scriptscriptstyle #2}$}
              \restraux{#1}{#2}}
              {\setbox1\hbox{${\scriptscriptstyle #1}_{\scriptscriptstyle #2}$}
              \restraux{#1}{#2}}}
\def\restraux#1#2{{#1\,\smash{\vrule height .8\ht1 depth .85\dp1}}_{\,#2}} 

%Quelques normes
\DeclarePairedDelimiter\abs{\lvert}{\rvert}
\DeclarePairedDelimiter\labs{\bigl\lvert}{\bigr\rvert}
\DeclarePairedDelimiter\norm{\lVert}{\rVert}
\DeclarePairedDelimiterXPP\normi[1]{}\lVert\rVert{_1}{\ifblank{#1}{\:\cdot\:}{#1}}
\DeclarePairedDelimiterXPP\normd[1]{}\lVert\rVert{_2}{\ifblank{#1}{\:\cdot\:}{#1}}
\DeclarePairedDelimiterXPP\normp[1]{}\lVert\rVert{_p}{\ifblank{#1}{\:\cdot\:}{#1}}
\DeclarePairedDelimiterXPP\normq[1]{}\lVert\rVert{_q}{\ifblank{#1}{\:\cdot\:}{#1}}
\DeclarePairedDelimiterXPP\norminf[1]{}\lVert\rVert{_\infty}{\ifblank{#1}{\:\cdot\:}{#1}}

%Bracket en angle
\DeclarePairedDelimiter\angl{\langle}{\rangle}

%Set, parenthèses
\newcommand{\set}[1]{\{ #1 \}}
\newcommand{\bigset}[1]{\bigl\{ #1 \bigr\}}
\newcommand{\Bigset}[1]{\Bigl\{ #1 \Bigr\}}
\newcommand{\bigpar}[1]{\bigl( #1 \bigr)}
\newcommand{\Bigpar}[1]{\Bigl( #1 \Bigr)}

%Opitmisation des opérateurs de comparaison
\DeclarePairedDelimiter\tio{o (}{)}
\DeclarePairedDelimiter\gro{O (}{)}


\newcommand{\dse}{développable en série entière }

\pagestyle{fancy}

\fancyhead[C]{}
\fancyhead[R]{}

\begin{document}

\underline{Source :} Gourdon - Analyse - page 261

\underline{Leçons :}
\begin{itemize}
\item 243 : Séries entières, propriétés de la somme. Exemples et applications.
\item 218 : Formules de Taylor. Exemples et applications.
\end{itemize}

\begin{thm}
Soient $a>0$ et $f : ]-a, a[ \to \R$ une fonction de classe $\Cinf$. On suppose que :
\[ \forall k \in \N, \forall x \in ]-a, a[, \quad f^{(2k)}(x) \geq 0 . \]
Alors $f$ est développable en série entière sur$]-a,a[$.
\end{thm}

\begin{proof}
Soit $b \in ]0,a[$, montrons que $f$ est \dse sur $]-b, b[$.

\begin{enumerate}[label=(\roman*)]

\item
On considère la fonction $F$ définie sur $]-a,a[$ par $F(x) = f(x) + f(-x)$. 
On peut appliquer la formule de Taylor avec reste intégral à l'ordre $p +1\in \N$ à $F$, on trouve, pour $x \in [0,b]$ :
\[ F(x) = \sum_{k=0}^p \dfrac{F^{k}(0)}{k!} x^k + \int_0^b \dfrac{F^{p+1}(t)}{p!} (x-t)^p dt  , \]
et la fonction $F$ étant paire on a, $F^{(k)}(0) = 0$ pour $k$ impair. Ainsi en écrivant $p=2n+2$ on trouve :
\[ F(x) = F(0) + \dfrac{F^{2}(0)}{2!} x^2 + \dots + \dfrac{F^{2n}(0)}{(2n)!} x^{2n} + R_n(x), \qquad 
R_n(x) = \int_0^x \dfrac{F^{2n+2}(t)}{(2n+1)!} (x-t)^{2n+1} . \]
On remarque que pour tout $k \in \N, F^{(2k)}(0) = 2 f^{(2k)}(0) \geq 0$ donc $0 \leq R_n(b) \leq F(b)$ pour tout $n \in \N$.
Remarquons aussi que :
\begin{align*}
R_n(x) & = \int_0^x \left ( \dfrac{x-t}{b-t} \right )^{2n+1} \dfrac{(b-t)^{2n+1}}{(2n+1)!} F^{2n+2}(t) dt \\
& \leq \left ( \dfrac{x}{b} \right )^{2n+1} \int_0^x \dfrac{(b-t)^{2n+1}}{(2n+1)!} F^{2n+2}(t) dt \quad 
\text{car } t \mapsto \dfrac{x-t}{b-t} \text{ décroissante sur } [0,x] \\
& \leq \left ( \dfrac{x}{b} \right )^{2n+1} R_n(b) \\
& \leq \left ( \dfrac{x}{b} \right )^{2n+1} F(b)
\end{align*}
Ainsi, pour $x \in [0,b[$, on a $\lim_{n \to \infty} R_n(x) = 0$, soit, en passant à la limite dans la formule de Taylor :
\[ \forall x \in [0,b[, \quad F(x) = \sum_{n=0}^{\infty} \dfrac{F^{n}(0)}{n!} x^n , \]
et, par parité, on peut étendre ce résultat sur $]-b,b[$. 

\item
On veut maintenant démontrer le résultat pour $f$. On fixe $x \in ]-b, b[$, et pour $n \in \N$, on écrit :
\[ f(x) = f(0) + \dots + \dfrac{x^{2n+1}}{(2n+1)!} f^{(2n+1)}(0) + r_n(x) \qquad
 r_n(x) = \int_0^x \dfrac{(x-t)^{2n+1}}{(2n+1)!} f^{(2n+2)}(t) dt . \]
Comme $0 \leq  f^{(2n+2)}(t) \leq  f^{(2n+2)}(t) +  f^{(2n+2)}(-t) =  F^{(2n+2)}(t)$.
Donc, en prolongeant $R_n$ sur $]-b, b[$ par parité, on obtient $0 \leq r_n(x) \leq R_n(x)$ sur $]-b, b[$ 
et donc $\lim_{n \to \infty} r_n(x) = 0$. En notant :
\[ \forall p \in \N, \quad S_p(x) = \sum_{k=0}^p \dfrac{x^k}{k!} f^{(k)}(0), \]
on a montré que :
\begin{equation} \lim_{n \to \infty} S_{2n+1}(x) = f(x) \label{eq:2} . \end{equation}
Or, pour $n \geq 1$ :
\[ S_{2n}(x) - S_{2n-1}(x) = \dfrac{x^{2n}}{(2n)!} f^{(2n)}(0) = \dfrac12 \dfrac{x^{2n}}{(2n)!} F^{(2n)}(0) \to 0 , \]
comme terme général d'une série convergente, et donc :
\[ \lim_{n \to \infty} S_{2n}(x) - S_{2n-1}(x) = 0 , \]
d'où on déduit de (\ref{eq:2}) que $\lim_{n \to \infty} S_n(x) = f(x)$ et donc que $f$ est \dse sur $]-b,b[$.
Ceci étant valable sur tout ouvert $]-b, b[ \subset ]-a,a[$, c'est vrai sur $]-a,a[$ tout entier.

\end{enumerate}
\end{proof}

\begin{appl}
On considère la fonction :
\[ \begin{array}{c c c} 
f: ] -\pi/2, \pi/2[ & \to & \R \\
x & \mapsto & \tan{x} 
\end{array} \]
La fonction $f'$ est définie sur ce même intervalle par $f'(x) = \dfrac1{\cos^2{x}}$ et on vérifie que les dérivées d'ordre paires succesives
de $f'$ sont des fractions rationnelles de monômes de degrés pairs évaluées en $\cos{x}$ et en $\sin{x}$.
Ainsi, $f'$ satistfait aux conditions du théorème de Bernstein, on conclut donc qu'elle est \dse sur $]-\pi/2, \pi/2[$
et par suite que la fonction $x \mapsto \tan{x}$ est \dse sur ce même intervalle.
\end{appl}

\end{document}